\section*{A sharp constant-child exclusion for critical FEI sources}

All Fourier entropies are in bits.  For a nonconstant Boolean function $f$,
write $R(f)=H(f)/I(f)$.  Call $f$ restriction-critical if every nonconstant
physical restriction $g$ of $f$ satisfies $R(g)\le R(f)$.

\subsection*{Source identity}

Let one physical coordinate $x$ split $K$ into a constant child $c\in
\{-1,1\}$ and a Boolean child $D$.  Multiplication by $c$ does not change
any squared Fourier coefficient, so assume $c=1$ and put
$p=\Pr[D=-1]$.  The independently reviewed scalar packet identity and the
sharp bias calculation give
\begin{align}
 H(K)&=\frac12H(D)+p\,c(p), &
 I(K)&=\frac12I(D)+p,                                      \tag{1}\\
 c(p)&\le C_*, &
 4.872321<C_*&<4.872322.                                   \tag{2}
\end{align}
Here
\[
c(p)=\frac1p\left[
2(1-p)^2\log_2\frac1{1-p}+2p^2\log_2\frac1p+4p(1-p)
-(1-2p)^2\log_2\frac1{|1-2p|}\right],                     \tag{3}
\]
with continuous endpoint conventions.  Equation (1) retains the full child
entropy and influence.  There is no independence assumption beyond the one
fresh physical selector already present in the restriction split.

For completeness, (2) is an all-dimensional one-variable statement.  With
$A(p)=(\ln2)p c(p)$, the sign of $c'(p)$ is the sign of
\[
G(p)=2(1-p^2)\ln(1-p)-2p^2\ln p-4p^2\ln2
 -(1-4p^2)\ln|1-2p|,                                      \tag{4}
\]
and
\[
G'(p)=4p\ln\frac{(1-2p)^2}{4p(1-p)}.                       \tag{5}
\]
The derivative changes sign only at
$\alpha=(1-1/\sqrt2)/2$ and $\beta=(1+1/\sqrt2)/2$.
Together with $G(0+)=0$, $G(1/2)=-2\ln2$, and
$G(1-)=-4\ln2$, this gives the unique maximizer.  The source-bound rational
interval certificate brackets it in $(0.22944127,0.22944128)$ and proves the
upper bound in (2).

\subsection*{Cross-route substitution}

Assume first that $D$ is nonconstant.  Put $R=R(K)$.  Restriction-criticality
at the physical child gives
\[
 H(D)\le R I(D).                                             \tag{6}
\]
Substituting (1)--(2) and (6) into the defining equality
$R I(K)=H(K)$ yields
\[
 R\left(\frac12 I(D)+p\right)
 =\frac12H(D)+pc(p)
 \le \frac R2 I(D)+C_*p.                                    \tag{7}
\]
Since $D$ is not the constant $1$, $p>0$.  Cancelling the child term and
dividing by $p$ gives
\[
 \boxed{R(K)\le C_*<4.872322.}                              \tag{8}
\]
If both children are constant, relevance makes $K$ a signed dictator and
$H(K)=0$, so (8) is again true.

Thus every restriction-critical Boolean source with ratio at least
$4.872322$ has no relevant coordinate with a constant child.  Combining this
with the critical-prime descent reduces any hypothetical source above that
threshold to a tensor-prime source whose every codimension-one physical face
is nonconstant.

\subsection*{Debt paid and debt remaining}

The earlier critical-pruning route asked for a universal bound on
$\min_i\Delta_iH/I_i$.  The substitution pays that obligation completely on
the constant-child stratum: such a coordinate cannot occur in a critical
source above $C_*$.  It also excludes every sparse-activation outer control
from a critical endpoint above $C_*$, while retaining its active child's full
entropy and influence.

The exact remaining class has two nonconstant children at every relevant
coordinate.  Addressing, bent splices, shared-variable interference, and
general near-tensor surgeries may lie in that class.  No estimate for them is
claimed.  In particular, (8) is a structural restriction on near-extremizers,
not a proof that full FEI has constant $C_*$.

